\scanpage{3}\mspage{1}
\deleted{\unsure{Les}} Retour \deleted{\unsure{sur|aux}}\unsure{sur} les catégories de fractions $M\Sigma^{-1}$.

1. \uline{Fonctorialité du} $\mathrm{Ch}^{\Sigma}_{\infty}(a,b)$ \uline{en} $a,b$ \uline{dans} $M\Sigma^{-1}$

On se donne une cat. $M$ avec ens. admissible de flèches $\Sigma\subset\mathrm{Fl}(M)$, d'où (en prenant p.ex. $\underline t_{\infty}=\bigcup\underline t_n$), avec $\underline t_n=\deleted{\unsure{\xi}}(\underbrace{\Delta_+\Delta_-}_{\tau_0})^{2n}$, et $\underset{\substack{\Vert\\\tau_0^{2n+2}}}{t_{n+1}}\unsure{\longrightarrow|\text{---}}\underset{\substack{\Vert\\\tau_0^{2n}}}{t_n}$ défini \deleted{par les} \deleted{\unsure{composé}} $p\unsure{*|\times}\mathrm{id}_{\tau_0^{2n}}\unsure{\times|*}p$, où $p:\tau_0\to1$ la dégénérescence. On a alors une théorie des $\mathrm{Ch}_{\Sigma;\infty}(M,a;b)$, que je note simplement $\mathrm{Ch}_{\deleted{\unsure{M}}\infty}(a,b)$\inserted{\unsure{en} \unsure{tant|vus} \unsure{qu'}objets de $\mathrm{Hot}$}. Des \deleted{\unsure{lois}} accouplements

\[
\mathrm{Ch}_{\deleted{\unsure{M}}\infty}(a,b)\times\mathrm{Ch}_{\deleted{\unsure{M}}\infty}(b,c)\longrightarrow\mathrm{Ch}_{\deleted{\unsure{M}}\infty}(a,c)\qquad\text{dans }\mathrm{Hot}
\]
avec leurs \unsure{liens} ($=$ admettant) \uline{associativité} \deleted{\unsure{en}} dans $\mathrm{Hot}$. \deleted{\unsure{Enfin, on a}} Cela définit donc des accouplements sur les $\pi_0(\mathrm{Ch}_{\deleted{\unsure{M}}\infty}(a,b))$\inserted{$\overset{\text{déf}}{=}\pi_0(a,b)$}, qui \unsure{n'est|n'est rien} autre que la composition des \deleted{flèches} \deleted{dans} $M\Sigma^{-1}$. D'autre part on a pour $a,b\in\mathrm{Ob}M$

\[
\mathrm{Hom}_M(a,b)\longrightarrow\pi_0(\mathrm{Ch}_{\infty}(a,b))=\pi_0(a,b)
\]
\unsure{compatible|compatibles} avec les compositions. Mais on peut \unsure{l'ignorer} pour l'instant, et regarder $\mathrm{Ob}M$, \inserted{ou simplement $M_0$,} avec la seule structure
\[
\left\{\begin{aligned}
M_0\times M_0&\longrightarrow\mathrm{Hot}&(a,b)&\unsure{\mapsto|\vdash}\mathrm{Ch}_{\infty}(a,b)\\
M_0\times M_0\times M_0&\longrightarrow\mathrm{Fl}\,\mathrm{Hot}&(a,\underline b,c)&\mapsto\bigl(\mathrm{Ch}_{\infty}(a,b)\times\mathrm{Ch}_{\infty}(b,c)\\
&&&\hspace{2em}\underset{\text{fl. dans Hot}}{\unsure{\twoheadrightarrow|\longrightarrow}}\mathrm{Ch}_{\infty}(a,c)\bigr)
\end{aligned}\right.
\]
avec \inserted{\unsure{condition}} associativité, et aussi une condition d'unités : $\mathrm{Ch}_{\infty}(a,a)$ est \unsure{muni|munie} d'un
\[
1_a\in\pi_0\,\mathrm{Ch}_{\infty}(a,a)=\pi_0(a,a),\quad\deleted{\unsure{\text{pour}}}
\]
\[
\deleted{\unsure{\pi_0(a,b)=\pi_0\mathrm{Ch}_0(a,b)}}
\]
jouant le rôle d'une ``identité bilatère'', dans le sens suivant. Notons que les accouplements
\[
\begin{aligned}
\mathrm{Ch}_{\infty}(a,b)\times\mathrm{Ch}_{\infty}(b,c)&\longrightarrow\mathrm{Ch}_{\infty}(a,c)\\
P\times Q&\longrightarrow R
\end{aligned}
\]
% archive sheet number: 2

